\section{Introduction}
\label{sec:intro}

Fluid--structure interaction (FSI) problems in which a soft solid undergoes large
deformation inside a flow---valves, swimmers, capsules, gels and polymeric
particles---are numerically demanding because the solid and fluid share a moving
interface across which stress and velocity must remain continuous. The dominant
strategies treat the two phases with different kinematics. Arbitrary
Lagrangian--Eulerian (ALE) methods deform a body-fitted mesh with the interface
\cite{Donea1982,Hu2001}, which is accurate but requires remeshing once the
deformation becomes severe. Immersed and fictitious-domain methods keep the fluid
grid fixed and represent the solid through forcing or Lagrange multipliers
\cite{Peskin2002,Zhang2004}, avoiding remeshing at the cost of interface sharpness.
Partitioned couplings that iterate between separate fluid and solid solvers are
modular but suffer the added-mass instability when the densities are comparable, so
that strong coupling or careful relaxation is required for stability
\cite{Causin2005,Hou2012}. These difficulties have motivated a third route: fully
\emph{Eulerian}, one-fluid formulations that discretize solid and fluid on a single
fixed grid, sharing one velocity field and one pressure solve
\cite{Sugiyama2011,Ii2012,Cottet2008,Wick2013}. The interface is captured rather than
tracked, the added-mass coupling is resolved monolithically inside the momentum
balance, and no remeshing is ever needed; the price is that the finite-strain solid
response, naturally posed on a material reference, must be reconstructed in the
laboratory frame.

The \emph{reference map technique} (RMT) supplies exactly this reconstruction. It
transports the inverse motion---the reference map $\xiv(\xv,t)$ carrying a current
point back to its material label---as an Eulerian field, from which the deformation
gradient $\Ften=(\grad\xiv)^{-1}$ and hence the elastic stress follow pointwise
\cite{Kamrin2012,Valkov2015}. Since its introduction the RMT has grown into a robust
incompressible FSI solver with level-set interface capture and stress blending
\cite{Rycroft2020}, and into a conservative, non-dissipative finite-difference
formulation \cite{Jain2019}; it belongs to the same Eulerian-solid lineage as the
full-Eulerian schemes of Sugiyama \emph{et al.}\ \cite{Sugiyama2011,Sugiyama2010} and
the earlier continuum description of viscoelastic particles by Liu and Walkington
\cite{Liu2001}. The appeal is a single set of conservation laws on one staggered
grid; the open questions concern \emph{what} constitutive branch the solid carries and
\emph{how} the coupled system is advanced in time.

Both questions are consequential. On the constitutive side, the established RMT
formulations carry a purely \emph{hyperelastic} solid \cite{Kamrin2012,Rycroft2020,Jain2019}.
A hyperelastic body sheared without bound stores unbounded energy, and its Eulerian
carrier---the reference map, or equivalently the elastic strain---folds under sustained
shear, so that reinitialization or filtering is used to treat the symptom. The
physically principled remedy is a \emph{relaxing} stress: a finite-strain viscoelastic
constitutive law bounds the stored elastic strain, and in the fluid limit removes the
folding sensitivity at its root rather than by numerical repair. Transporting a
viscoelastic stress stably is itself a classical difficulty---the high-Weissenberg
number problem, in which a naive update of the conformation tensor loses
positive-definiteness and the computation breaks down \cite{Oldroyd1950,Bird1987}.
The cure is to evolve a matrix logarithm of the conformation tensor, the
log-conformation representation of Fattal and Kupferman \cite{Fattal2004,Fattal2005},
subsequently recast in eigenframe form by Hulsen \emph{et al.}\ \cite{Hulsen2005} and
placed in a broader factorization framework by Balci \emph{et al.}\ \cite{Balci2011},
which guarantees a symmetric positive-definite conformation by construction. Bringing
this representation into the RMT ties the relaxing elastic strain of the solid to the
same stable transport that viscoelastic \emph{fluids} already enjoy, and decouples it
from the folding-prone reference map.

On the time-integration side, the same formulations advance momentum explicitly, so
the step is throttled simultaneously by the elastic-wave speed $c_s=\sqrt{\mu_s/\rho}$,
the viscous diffusion $\dx^2/\nu$, the advective CFL, and---once the solid relaxes---the
relaxation time $\tau$. Stiff solids and high-Weissenberg regimes are then effectively
out of reach. Implicit and monolithic Eulerian FSI does exist, but for
\emph{hyperelastic} solids: Dunne \cite{Dunne2006}, Richter \cite{Richter2013,Richter2017}
and Wick \cite{Wick2013} solve a coupled Newton system for velocity, pressure and a
solid deformation field on a fixed grid, tracking the interface with an ``initial point
set'' that is itself a reference-map-like field. Even there the deformation/interface is
advanced through a separate extension step outside the Jacobian, a treatment that is not
\emph{fully} implicit in the geometry. Sharp-interface implicit FSI for large-deformation
solids has likewise been developed by Kolahdouz \emph{et al.}\ \cite{Kolahdouz2023}, but
through a \emph{partitioned} Dirichlet--Neumann coupling on a separate Lagrangian solid
mesh rather than a monolithic Eulerian solve. In none of these is a finite-strain
\emph{viscoelastic} solid coupled to a fluid and advanced implicitly on a single grid.

This paper develops that method. Its contribution is the numerical formulation itself:
a semi-implicit, log-conformation Reference Map Technique for incompressible
viscoelastic FSI on a marker-and-cell (MAC) staggered grid, together with a family of
time integrators that remove the stiff step-size limits one at a time. The solid stress
is carried by a log-conformation tensor $\psib=\log\be$, positive-definite by
construction and, in the standard-linear-solid form, retaining an equilibrium
hyperelastic branch so the body still returns to shape while $\tau\!\to\!\infty$
recovers the neo-Hookean RMT exactly. The staggered discretization makes the pressure
projection exact---divergence is driven to machine zero---removing the
checkerboard and Rhie--Chow workarounds that a collocated arrangement requires
\cite{Harlow1965,Chorin1968,Brown2001}. On top of this spatial method we build four
time-integration schemes of increasing implicitness. An explicit scheme reproduces the
established RMT and fixes the reference. An implicit--explicit (IMEX) scheme treats the
viscous term by backward Euler through the pressure transforms and Krylov solvers, the
relaxation by an exact exponential update, and the elastic coupling by an
energy-conserving trapezoidal rule, lifting the viscous, relaxation and elastic-wave
limits in a fractional-step setting. A semi-Lagrangian variant additionally removes the
advective CFL restriction. Finally, a Jacobian-free Newton--Krylov (JFNK) scheme solves
the coupled velocity--pressure--reference-map system \emph{semi-monolithically},
preconditioned by one step of the split IMEX solver with the exact projection acting as
the pressure Schur complement, and treats advection and the elastic coupling implicitly
and consistently. Each scheme is verified to reproduce the explicit result before its
target limit is lifted, and we report a work--precision comparison across the four,
delineating the regimes---stiff solids, high Weissenberg number---where implicit
integration pays for itself. The consistent linearized neo-Hookean tangent underlying
the implicit coupling, and the physics-based preconditioning that makes the Newton
solve tractable, follow established practice for stiff nonlinear systems
\cite{Knoll2004,Kelley1995,Ascher1997,Reese1998,Simo1998}.

We are deliberate about scope. The IMEX and semi-Lagrangian schemes are fully
implemented and validated; the JFNK scheme is realized for the fluid and the
neo-Hookean elastic coupling, and its extension to the coupled log-conformation block
is identified as the remaining step. Every capability is reported at its true status.
The remainder of the paper is organized as follows.
\Cref{sec:governing,sec:rmt,sec:constitutive} set out the continuous one-fluid model,
the reference map, and the finite-strain viscoelastic constitutive law;
\cref{sec:discretization,sec:time} give the MAC discretization and the four
time-integration schemes with their algorithms;
and \cref{sec:verification,sec:validation,sec:results} present the verification against
analytic and manufactured limits, the FSI validation against the disc-in-vortex and
lid-driven-cavity benchmarks, and the work--precision comparison.
